\documentclass[11pt]{article}
\usepackage{ochamath}
\subjectcolor{2F5DA8}
\setsheet{線形代数}{対角化　演習}{https://university-math-crj.pages.dev/linear-algebra/diagonalization/}
\begin{document}
\makesheethead{解答}

\problem[対角化]
\begin{ans}
固有値は3と2。固有ベクトルはそれぞれ $(1,0)^{\mathsf T}$、$(-1,1)^{\mathsf T}$。\[P=\begin{pmatrix}1&-1\\0&1\end{pmatrix},\quad P^{-1}=\begin{pmatrix}1&1\\0&1\end{pmatrix},\quad D=\operatorname{diag}(3,2).\] $AP=PD$ を掛け算で確認できる。
\end{ans}
\point{列の順番と対角成分の順番を揃える。}

\problem[累乗]
\begin{ans}
$D^2=\operatorname{diag}(9,4)$ より \[A^2=PD^2P^{-1}=\begin{pmatrix}9&5\\0&4\end{pmatrix}.\] 直接の積では右上が $3\times1+1\times2=5$ となる。
\end{ans}
\point{対角成分の累乗だけを先に計算できる。}
\newpage

\problem[対角化できない例]
\begin{ans}
特性多項式は $(\lambda-2)^2$。$(J-2I)\boldsymbol v=\boldsymbol0$ から $v_2=0$ なので固有空間の次元は1。2本の独立な固有ベクトルが必要なのに1本分しかなく、対角化できない。
\end{ans}
\point{固有値の個数と固有ベクトルの独立な本数を区別する。}

\problem[状態遷移]
\begin{ans}
$M(2,3)^{\mathsf T}=(2,3)^{\mathsf T}$、$M(1,-1)^{\mathsf T}=0.5(1,-1)^{\mathsf T}$。固有値は1と0.5。初期値は $\frac15(2,3)^{\mathsf T}+\frac35(1,-1)^{\mathsf T}$ なので、1回後は $(0.7,0.3)^{\mathsf T}$、総和は1。定常分布は $(2/5,3/5)^{\mathsf T}$。
\end{ans}
\point{定常成分は固有値1、差の成分は0.5倍になる。}

\end{document}
